\chapter{Ketamine}
\index{Ketamine}

\section*{Definition and Overview}
Ketamine is a medicine used as an anaesthetic and pain reliever in hospitals, and in recent years it has been used at low doses to treat severe depression that has not responded to other treatments. It is also used illegally as a recreational drug, known as Special K, where it produces dissociation, hallucinations, and sedation. Recreational use carries significant risks, including bladder and kidney damage with chronic use, accidents and injuries while intoxicated, dependence, and the danger of overdose, especially when mixed with alcohol or other drugs. This chapter explains what ketamine is, its medical uses, the risks of recreational use, and where to get help.

\section*{Epidemiology}
Ketamine is a widely used medical anaesthetic, and its use in depression treatment has grown substantially, with ketamine and esketamine clinics operating across some countries. Recreational ketamine use has increased, particularly among young adults at parties and festivals, and it is one of the drugs most commonly implicated in hospital presentations for intoxication in some regions. Regular use is associated with serious bladder damage, and cases of ketamine-induced cystitis have risen sharply in countries with high use. Data on the full extent of recreational use are limited, but the harms are well documented.

\section*{Aetiology and Pathophysiology}
Ketamine works by blocking the NMDA receptor, a key receptor for the neurotransmitter glutamate, which disrupts normal brain signalling and produces anaesthesia at high doses and dissociation at low doses. In depression treatment, this blockade is thought to trigger changes in brain connections that relieve symptoms rapidly, sometimes within hours. At recreational doses, ketamine causes a trance-like state with distorted perception and out-of-body experiences. Chronic use damages the bladder lining, causing painful, frequent urination and, in severe cases, bladder shrinkage and kidney damage, and it also harms cognitive function, memory, and the liver and bile ducts. Dependence can develop, and withdrawal causes cravings, anxiety, and disturbed sleep.

\section*{Clinical Features}
\begin{itemize}
    \item Medical uses: anaesthesia, pain relief, and treatment-resistant depression
    \item Depression treatment: low-dose infusions or nasal sprays under medical supervision
    \item Recreational effects: dissociation, numbness, hallucinations, and euphoria
    \item Impaired coordination, balance, and judgement while intoxicated
    \item Slurred speech, confusion, and a blank stare
    \item Chronic use: bladder pain, frequent and painful urination, and blood in the urine
    \item Memory and concentration problems with regular use
    \item Dependence, with cravings and withdrawal
    \item Overdose: severe sedation, breathing depression, and unconsciousness, especially with alcohol
\end{itemize}

\section*{Differential Diagnosis}
The effects of ketamine use may be confused with other conditions.
\begin{itemize}
    \item \textbf{Other drug effects:} alcohol, GHB, and hallucinogens
    \item \textbf{Bladder conditions:} urinary tract infection, interstitial cystitis, and bladder cancer
    \item \textbf{Psychiatric conditions:} psychosis, dissociation, and anxiety
    \item \textbf{Neurological conditions:} seizures, stroke, and head injury
    \item \textbf{Liver and bile duct disease:} from chronic use
    \item \textbf{Depression:} which may prompt medical ketamine treatment
\end{itemize}

\section*{When to Seek Medical Care}
Anyone experiencing bladder symptoms, abdominal pain, memory problems, or mental health concerns related to ketamine use should seek medical help, and people using ketamine regularly should be assessed for dependence and harm. Urgent care is needed for severe intoxication, breathing difficulty, or injury while under the influence. Families concerned about a person's use should seek advice early.

\textbf{Contact your local emergency services for:}
\begin{itemize}
    \item Unconsciousness, difficulty breathing, or severe sedation after ketamine use
    \item Severe agitation, confusion, or hallucinations
    \item Chest pain, a racing heartbeat, or collapse
    \item Severe abdominal or bladder pain
    \item Injuries from falls or accidents while intoxicated
    \item Signs of a drug overdose mixed with alcohol or other drugs
\end{itemize}

\section*{Diagnosis and Investigations}
\begin{itemize}
    \item \textbf{History:} pattern of use, dose, and associated substances
    \item \textbf{Assessment of intoxication:} observation and vital signs
    \item \textbf{Urine tests:} for bladder damage and drug screening
    \item \textbf{Kidney and liver function tests:} for chronic effects
    \item \textbf{Imaging:} ultrasound or cystoscopy for bladder damage
    \item \textbf{Mental health assessment:} for depression, psychosis, and dependence
    \item \textbf{Referral:} to drug and alcohol services and specialists
\end{itemize}

\section*{Management}
\begin{itemize}
    \item \textbf{Emergency care:} breathing support and monitoring for overdose
    \item \textbf{Stopping use:} the key to reversing bladder and other damage
    \item \textbf{Treatment of bladder damage:} urology review and symptom management
    \item \textbf{Dependence treatment:} counselling, withdrawal support, and drug and alcohol programs
    \item \textbf{Mental health care:} treatment of depression and other conditions
    \item \textbf{Medical ketamine treatment:} for treatment-resistant depression, under specialist care
    \item \textbf{Harm reduction:} education about the risks of mixing drugs and driving
    \item \textbf{Support services:} counselling, peer support, and rehabilitation
\end{itemize}

\section*{Complications}
\begin{itemize}
    \item Ketamine bladder syndrome: pain, frequency, and damage to the bladder
    \item Kidney damage and, rarely, kidney failure
    \item Liver and bile duct damage
    \item Cognitive impairment and memory loss with chronic use
    \item Accidents, injuries, and risky behaviour while intoxicated
    \item Dependence and withdrawal
    \item Overdose and death, especially with alcohol or other depressants
    \item Psychological dependence and effects on work and relationships
\end{itemize}

\section*{Prognosis and Outcomes}
Early bladder damage from ketamine often improves when use stops, and stopping use is the most important step. Dependence responds to structured treatment, and most people recover with support. The medical use of ketamine for depression is generally safe and effective under supervision, with careful screening. The harms of recreational use are significant but largely preventable and reversible when use stops in time.

\section*{Prevention and Self-Care}
\begin{itemize}
    \item Be aware of the serious risks of recreational ketamine use
    \item Never mix ketamine with alcohol or other depressant drugs
    \item Do not drive or operate machinery after use
    \item Seek help early for bladder pain, frequency, or blood in the urine
    \item If using ketamine for depression, do so only under medical supervision
    \item Talk to a doctor about dependence or concerns about your use
    \item Support friends and family members to seek help
    \item Access drug and alcohol services for treatment and harm reduction
    \item Seek help for underlying depression or anxiety through proper channels
\end{itemize}

\section*{Sources and Further Reading}
The following reputable sources provide further information for healthcare students and patients seeking reliable, current advice about ketamine.
\begin{itemize}
    \item Alcohol and Drug Foundation. \textit{Ketamine information.} https://adf.org.au/
    \item Healthdirect Australia. \textit{Ketamine.} https://www.healthdirect.gov.au/
    \item Your Room (NSW Health). \textit{Ketamine facts.} https://yourroom.health.nsw.gov.au/
    \item Beyond Blue. \textit{Depression treatment information.} https://www.beyondblue.org.au/
    \item NHS. \textit{Ketamine.} https://www.nhs.uk/
    \item National Institute on Drug Abuse. \textit{Ketamine.} https://nida.nih.gov/
\end{itemize}
